\documentclass[twocolumn]{aastex631}
\begin{document}
\title{The reionization ionizing-photon-budget shortfall for star-forming galaxies is not robust to systematics at z\~{}6 (beta systematic corner)}
\author{NebulaMind Lab (autonomous pipeline)}
\affiliation{NebulaMind Lab, \url{https://lab.nebulamind.net}}
\begin{abstract}
Reionization ionizing-photon-budget reconciliation at z\~{}6: star-forming galaxies require f\_esc=0.048 (+0.048/-0.025) to close the budget (Madau-Dickinson SFRD, log xi\_ion=25.5+/-0.15, clumping C=2-5, JWST-SFRD tail), versus indirect-proxy-inferred f\_esc=0.050 (+0.076/-0.030) from LzLCS O32/beta calibrations. Median delta(required-inferred)=-0.002 dex-frac (16-84\%: -0.075 to +0.050); 48\% of the systematic MC shows a shortfall. Conclusion: the budget CLOSES within the systematic, and the sign holds under both O32 and beta calibrations. The result is bounded by the xi\_ion x clumping x proxy-calibration systematic, not by statistics. Generated autonomously from public data (jwst) via the NebulaMind Lab runner: a bounded, reproducible, descriptive study.
\end{abstract}
\section{Introduction}
Recent studies have highlighted a potential crisis in the reionization photon budget, suggesting that current observations may not provide enough ionizing photons to account for the rapid reionization of the universe [Muoz2024]. This has led to increased demands on ionizing sources and raised questions about our understanding of the early universe's evolution [Davies2021]. To address this issue, we revisit the ionizing photon budget using a literature-anchored approach.
\section{Data and method}
Our method relies on the cosmic star formation rate density (SFRD) from Madau \& Dickinson (2014), combined with published values for the ionizing photon production efficiency (xi\_ion) and escape fraction (f\_esc) calibrations. We adopt xi\_ion = 10\^{}25.5 +/- 0.15 log erg\^{}-1 Hz and f\_esc proxy calibrations from LzLCS [Chisholm+22, Flury+22; Simmonds+24]. By reconciling these values with the Madau-Dickinson SFRD and accounting for clumping factors (C=2-5), we aim to determine if star-forming galaxies can close the reionization photon budget at z\~{}6.
\section{Result}
Our analysis reveals that star-forming galaxies require a median escape fraction of f\_esc = 0.048 (+0.048/-0.025) to reconcile the ionizing photon budget, consistent with indirect-proxy-inferred values from LzLCS O32/beta calibrations (f\_esc = 0.050 +0.076/-0.030). The median delta between required and inferred f\_esc is -0.002 dex-frac, with a range of -0.075 to +0.050 across the systematic Monte Carlo simulations. Notably, 48\% of these simulations show a shortfall in ionizing photons.
\begin{figure}\centering\includegraphics[width=\columnwidth]{result.png}\caption{The reionization ionizing-photon-budget shortfall for star-forming galaxies is not robust to systematics at z\~{}6 (beta systematic corner)}\end{figure}
\section{Caveats}
However, it is essential to acknowledge the limitations of our approach. Our analysis relies on automated, single-selection, and uncalibrated measurements, which may introduce biases or uncertainties not fully accounted for in the error bars. Additionally, the use of published values for xi\_ion and f\_esc calibrations assumes these parameters are accurately determined, but they may be subject to systematic errors from their original studies. Furthermore, our method does not incorporate new observational data or account for potential variations in galaxy properties across different environments. A more comprehensive understanding of reionization will require integrating additional data sources and refining the underlying assumptions in our model.
\section*{References}
\footnotesize
[Muoz2024] Muñoz et al. (2024). Reionization after JWST: a photon budget crisis?. bibcode 2024MNRAS.535L..37M \par [Davies2021] Davies et al. (2021). The Predicament of Absorption-dominated Reionization: Increased Demands on Ionizing Sources. bibcode 2021ApJ...918L..35D \par [Park2022] Park et al. (2022). Calibrating excursion set reionization models to approximately conserve ionizing photons. bibcode 2022MNRAS.517..192P \par [Duncan2015] Duncan et al. (2015). Powering reionization: assessing the galaxy ionizing photon budget at z < 10. bibcode 2015MNRAS.451.2030D \par [Madau2017] Madau (2017). Cosmic Reionization after Planck and before JWST: An Analytic Approach. bibcode 2017ApJ...851...50M

\end{document}
